\documentclass[11pt]{article}
\usepackage[margin=1in]{geometry}
\usepackage{amsmath,amssymb,amsthm}
\usepackage{hyperref}

\newtheorem{assumption}{Assumption}
\newtheorem{lemma}{Lemma}
\newtheorem{proposition}{Proposition}
\newtheorem{theorem}{Theorem}
\newtheorem{corollary}{Corollary}
\newcommand{\R}{\mathbb{R}}
\newcommand{\supp}{\operatorname{supp}}

\title{Finite-Tolerance Refinement with Inexact Restricted KL-UOT Solves\\\large Working Theory Note}
\author{}
\date{October 12, 2026}

\begin{document}
\maketitle

\begin{abstract}
This note studies active-support refinement when restricted KL-UOT subproblems are solved only approximately.
The main observation is specific to unbalanced transport: increasing one missing transport edge is always primal feasible because the marginals are penalized rather than imposed as hard constraints.
At any coupling with positive transported marginals, a positive primal-consistent dual score therefore gives an explicit one-edge objective decrease.
We use this decrease to prove finite termination to a prescribed full-space tolerance for a safeguarded inexact-oracle refinement scheme.
The theorem requires the updated support to preserve the current positive flow and the inexact restricted solve to retain a fixed fraction of the certified one-edge decrease.
For nonnegative costs and node-covering restricted supports, an explicit curvature bound follows from the objective sublevel set and the restricted projected-gradient residual; a curvature-scaled stopping rule gives an alternative.
These conditions identify a small certified variant of HierUOT rather than proving the current production refinement loop unchanged.
\end{abstract}

\section{Scope and relation to prior work}

We consider one hierarchy level of KL-unbalanced optimal transport on the complete edge set $\mathcal E=[m]\times[n]$.
The equivalent-transformation mechanism of Liu et al.~\cite{liu2025etm} determines the optimal transported marginals and then reduces KL-UOT to balanced OT.
Consequently, existence of a forest-supported optimal coupling follows from standard extreme-point theory for the resulting balanced transport problem.
We use that fact only as background and do not claim it as a contribution here.

The target of this note is different.
We ask whether active-support refinement can terminate at a finite tolerance when each restricted KL-UOT call returns an approximate primal--dual point rather than an exact restricted optimum.
The result below exploits a feasible one-edge direction that has no analogue in balanced OT.

\section{Primal-consistent scores}

For $X\in\R_+^{m\times n}$, let $r=X\mathbf 1_n$ and $q=X^\top\mathbf 1_m$.
The full KL-UOT objective is
$$\Phi(X)=\langle C,X\rangle+\rho_a\sum_i\left(r_i\log\frac{r_i}{a_i}-r_i+a_i\right)+\rho_b\sum_j\left(q_j\log\frac{q_j}{b_j}-q_j+b_j\right),$$
where $a_i,b_j,\rho_a,\rho_b$ are positive and all costs are finite.

Suppose first that every entry of $r$ and $q$ is positive.
Define the primal-consistent potentials and scores by
$$f_i^X=-\rho_a\log\frac{r_i}{a_i},\qquad g_j^X=-\rho_b\log\frac{q_j}{b_j},\qquad \sigma_{ij}^X=f_i^X+g_j^X-C_{ij}.$$
The coordinate derivative of the primal objective is
$$\nabla_{ij}\Phi(X)=C_{ij}+\rho_a\log\frac{r_i}{a_i}+\rho_b\log\frac{q_j}{b_j}=-\sigma_{ij}^X.$$

For a fixed stepsize parameter $\gamma>0$ and an active support $\mathcal A\subseteq\mathcal E$, define the restricted projected-gradient residual
$$R_{\mathcal A}^{\gamma}(X)=\frac{1}{\gamma}\left\|X_{\mathcal A}-\left[X_{\mathcal A}-\gamma\nabla_{\mathcal A}\Phi(X)\right]_+\right\|_{\infty}$$
and the inactive full-space violation
$$V_{\mathcal A}(X)=\max_{e\in\mathcal E\setminus\mathcal A}(\sigma_e^X)_+.$$

\begin{proposition}[Global KKT residual]\label{prop:global-kkt}
Extend a coupling supported on $\mathcal A$ by zero outside $\mathcal A$.
Then
$$R_{\mathcal E}^{\gamma}(X)=\max\{R_{\mathcal A}^{\gamma}(X),V_{\mathcal A}(X)\}.$$
In particular, $R_{\mathcal A}^{\gamma}(X)=V_{\mathcal A}(X)=0$ if and only if $X$ is a global KL-UOT optimum.
\end{proposition}

\begin{proof}
For every inactive coordinate, $X_e=0$ and its projected-gradient residual is
$$\frac{1}{\gamma}\left|0-[0-\gamma\nabla_e\Phi(X)]_+\right|=(-\nabla_e\Phi(X))_+=(\sigma_e^X)_+.$$
Splitting the maximum over active and inactive coordinates proves the identity.
The projected-gradient residual vanishes exactly when the convex first-order optimality condition holds.
\end{proof}

Proposition~\ref{prop:global-kkt} shows that primal-consistent scores provide a complete finite-tolerance KKT certificate when combined with a restricted projected-gradient residual.
No primal objective-gap conclusion is implied without an additional error bound.

\section{One-edge descent}

\begin{lemma}[Certified one-edge decrease]\label{lem:one-edge}
Let $X\geq0$ have positive transported marginals and let $e=(i,j)$ be any edge.
Define
$$h_e(X)=\frac{\rho_a}{r_i}+\frac{\rho_b}{q_j}.$$
If $\sigma_e^X>0$, then the step $t_e=\sigma_e^X/h_e(X)$ satisfies
$$\Phi(X+t_eE_e)\leq\Phi(X)-\frac{(\sigma_e^X)^2}{2h_e(X)},$$
where $E_e$ contains one at edge $e$ and zero elsewhere.
\end{lemma}

\begin{proof}
Let $\psi(t)=\Phi(X+tE_e)-\Phi(X)$ for $t\geq0$.
Direct differentiation gives
$$\psi'(0)=-\sigma_e^X,\qquad \psi''(t)=\frac{\rho_a}{r_i+t}+\frac{\rho_b}{q_j+t}\leq h_e(X).$$
The integral form of Taylor's theorem yields
$$\psi(t)=-\sigma_e^Xt+\int_0^t(t-s)\psi''(s)\,ds\leq-\sigma_e^Xt+\frac{h_e(X)}{2}t^2.$$
Substitution of $t=t_e$ proves the claim.
\end{proof}

The lemma is global and does not assume that $X$ is a restricted optimum or even an approximate restricted stationary point.
Its denominator is the same local curvature quantity that appears in curvature-aware pricing.
The score $(\sigma_e^X)^2/(2h_e(X))$ is therefore a guaranteed primal decrease, although ranking edges by this bound need not improve practical refinement time.

If a transported marginal is zero, the corresponding KL directional quotient tends to $-\infty$ along any incident edge.
Such a node also admits a strict one-edge decrease, but the finite bound in Lemma~\ref{lem:one-edge} is unavailable until the marginal becomes positive.

\subsection{Using approximate dual potentials}

The production full-space scan uses PDUOT potentials $(f,g)$ rather than the primal-consistent pair $(f^X,g^X)$.
Define the nodewise response mismatches
$$\eta_i^a=|f_i-f_i^X|,\qquad \eta_j^b=|g_j-g_j^X|.$$
For the scanned score $\sigma_{ij}=f_i+g_j-C_{ij}$, we have the two-sided enclosure
$$\sigma_{ij}-\eta_i^a-\eta_j^b\leq\sigma_{ij}^X\leq\sigma_{ij}+\eta_i^a+\eta_j^b.$$

\begin{corollary}[Dual-score correction]\label{cor:dual-correction}
Let $\delta_{ij}=(\sigma_{ij}-\eta_i^a-\eta_j^b)_+$.
If $\delta_{ij}>0$, the step $t=\delta_{ij}/h_{ij}(X)$ satisfies
$$\Phi(X+tE_{ij})\leq\Phi(X)-\frac{\delta_{ij}^2}{2h_{ij}(X)}.$$
If $\max_{(i,j)\notin\mathcal A}(\sigma_{ij}+\eta_i^a+\eta_j^b)_+\leq\varepsilon_f$, then $V_{\mathcal A}(X)\leq\varepsilon_f$.
\end{corollary}

\begin{proof}
The lower enclosure gives $\sigma_{ij}^X\geq\delta_{ij}$.
Applying the quadratic upper bound in the proof of Lemma~\ref{lem:one-edge} at $t=\delta_{ij}/h_{ij}(X)$ gives the first conclusion.
The upper enclosure gives the second conclusion.
\end{proof}

Let $\widehat r_i=a_i e^{-f_i/\rho_a}$ and $\widehat q_j=b_j e^{-g_j/\rho_b}$ be the dual-implied marginal responses.
The mismatches can be computed exactly as
$$\eta_i^a=\rho_a\left|\log\frac{r_i}{\widehat r_i}\right|,\qquad \eta_j^b=\rho_b\left|\log\frac{q_j}{\widehat q_j}\right|.$$
If $r_i,\widehat r_i\geq\underline m_a>0$, then the mean-value theorem gives $\eta_i^a\leq\rho_a|r_i-\widehat r_i|/\underline m_a$, with the analogous target bound.
Thus the existing marginal-response residual controls score uncertainty only when the relevant primal and dual-implied marginals have a positive lower bound.

The lower and upper enclosures play different roles.
A positive lower enclosure certifies a descent direction, whereas a small upper enclosure certifies full-space stationarity.
If neither condition holds, the restricted oracle must be tightened or the full scan can be repeated directly with $(f^X,g^X)$.

\section{Safeguarded inexact refinement}

We now state a finite-tolerance result for a certified refinement variant.
At round $k$, an inexact restricted oracle returns $X^k$ supported on $\mathcal A_k$ with positive marginals and $R_{\mathcal A_k}^{\gamma}(X^k)\leq\varepsilon_r$.
If $V_{\mathcal A_k}(X^k)\leq\varepsilon_f$, the level terminates.
Otherwise, the update selects an inactive edge $e_k$ satisfying $\sigma_{e_k}^{X^k}>\varepsilon_f$ and forms
$$t_k=\frac{\sigma_{e_k}^{X^k}}{h_{e_k}(X^k)},\qquad Y^k=X^k+t_kE_{e_k}.$$
The next support contains $\supp(X^k)\cup\{e_k\}$, so $Y^k$ remains feasible after pruning.
The next inexact restricted solve is required to return $X^{k+1}$ satisfying
$$R_{\mathcal A_{k+1}}^{\gamma}(X^{k+1})\leq\varepsilon_r,\qquad \Phi(X^{k+1})\leq\Phi(Y^k)+\xi_k.$$
The second condition is an objective safeguard rather than an objective-error certificate for the unrestricted problem.

\begin{theorem}[Finite-tolerance refinement with an inexact oracle]\label{thm:inexact-refinement}
Assume that $\Phi$ is bounded below by $\Phi_{\inf}$ and that every nonterminal round satisfies
$$h_{e_k}(X^k)\leq H,\qquad 0\leq\xi_k\leq\theta\frac{(\sigma_{e_k}^{X^k})^2}{2h_{e_k}(X^k)}$$
for constants $H<\infty$ and $\theta\in[0,1)$.
Then the safeguarded refinement scheme terminates after at most
$$N_{\mathrm{ref}}\leq\left\lceil\frac{2H(\Phi(X^0)-\Phi_{\inf})}{(1-\theta)\varepsilon_f^2}\right\rceil$$
nonterminal rounds.
Its output satisfies
$$R_{\mathcal E}^{\gamma}(X)\leq\max\{\varepsilon_r,\varepsilon_f\}.$$
\end{theorem}

\begin{proof}
Lemma~\ref{lem:one-edge} and the inexact-oracle safeguard give
$$\Phi(X^{k+1})\leq\Phi(X^k)-(1-\theta)\frac{(\sigma_{e_k}^{X^k})^2}{2h_{e_k}(X^k)}\leq\Phi(X^k)-\frac{(1-\theta)\varepsilon_f^2}{2H}.$$
If the number of nonterminal rounds exceeded the displayed bound, telescoping this inequality would force an iterate below $\Phi_{\inf}$.
The scheme must therefore terminate in finitely many rounds.
At termination, the restricted oracle gives $R_{\mathcal A}^{\gamma}(X)\leq\varepsilon_r$ and the full scan gives $V_{\mathcal A}(X)\leq\varepsilon_f$.
Proposition~\ref{prop:global-kkt} proves the final estimate.
\end{proof}

For the squared-Euclidean costs used by HierUOT, $C\geq0$ and every generalized KL term is nonnegative, so one may take $\Phi_{\inf}=0$.

The next proposition removes the separate curvature assumption when the restricted supports cover all nodes.
For $s>0$ and $B\geq0$, let $M(s,B)$ be the unique number in $[s,\infty)$ satisfying
$$M(s,B)\log\frac{M(s,B)}{s}-M(s,B)+s=B.$$

\begin{proposition}[Explicit marginal lower bound]\label{prop:marginal-lower-bound}
Assume $0\leq C_{ij}\leq C_{\max}$, $\Phi(X)\leq B$, the support $\mathcal A$ covers every source and target node, and $R_{\mathcal A}^{\gamma}(X)\leq\varepsilon_r$.
Define
$$L_a(B)=\max_i\log\frac{M(a_i,B/\rho_a)}{a_i},\qquad L_b(B)=\max_j\log\frac{M(b_j,B/\rho_b)}{b_j},$$
$$\underline r_i(B)=a_i\exp\left(\frac{-\varepsilon_r-C_{\max}-\rho_bL_b(B)}{\rho_a}\right),\qquad \underline q_j(B)=b_j\exp\left(\frac{-\varepsilon_r-C_{\max}-\rho_aL_a(B)}{\rho_b}\right).$$
Then $r_i\geq\underline r_i(B)$ and $q_j\geq\underline q_j(B)$ for every node.
Consequently,
$$h_{ij}(X)\leq H(B):=\frac{\rho_a}{\min_i\underline r_i(B)}+\frac{\rho_b}{\min_j\underline q_j(B)}.$$
\end{proposition}

\begin{proof}
Nonnegativity of the cost and KL terms implies
$$\rho_a\left(r_i\log\frac{r_i}{a_i}-r_i+a_i\right)\leq B,\qquad \rho_b\left(q_j\log\frac{q_j}{b_j}-q_j+b_j\right)\leq B.$$
Monotonicity of the generalized KL divergence above its minimizer gives $r_i\leq M(a_i,B/\rho_a)$ and $q_j\leq M(b_j,B/\rho_b)$.
The projected-gradient residual bound implies $\nabla_e\Phi(X)\geq-\varepsilon_r$ for every $e\in\mathcal A$.
Indeed, this is immediate when the gradient is nonnegative, whereas a negative gradient makes the projection inactive and its absolute value equals the coordinate residual.
For each source node $i$, choose an incident active edge $(i,j)$ and rearrange its gradient inequality to obtain
$$\rho_a\log\frac{r_i}{a_i}\geq-\varepsilon_r-C_{ij}-\rho_b\log\frac{q_j}{b_j}\geq-\varepsilon_r-C_{\max}-\rho_bL_b(B).$$
Exponentiation proves the source bound, and the target bound is symmetric.
Substitution in the definition of $h_{ij}$ proves the last claim.
\end{proof}

\begin{corollary}[Data-dependent refinement bound]\label{cor:data-dependent-bound}
Suppose the hypotheses of Theorem~\ref{thm:inexact-refinement} hold except for the assumed constant $H$, every restricted support covers all nodes, and $C\geq0$.
Then the theorem remains valid with $H=H(\Phi(X^0))$ from Proposition~\ref{prop:marginal-lower-bound}.
\end{corollary}

\begin{proof}
Every nonterminal update decreases the objective, so $\Phi(X^k)\leq\Phi(X^0)$.
Proposition~\ref{prop:marginal-lower-bound} supplies the required uniform curvature bound at every restricted-oracle output.
\end{proof}

The explicit constant is generally too conservative to predict the observed number of refinement rounds.
Its role is qualitative: it closes the finite-termination argument using verifiable problem data and oracle conditions rather than a separate lower-marginal assumption.

\subsection{A curvature-scaled alternative}

Define the maximum certified one-edge gain
$$G_{\mathcal A}(X)=\max_{e\notin\mathcal A}\frac{((\sigma_e^X)_+)^2}{2h_e(X)}.$$

\begin{corollary}[Gain-based termination]\label{cor:gain-termination}
Suppose a nonterminal round is executed only when $G_{\mathcal A_k}(X^k)>\varepsilon_G$, the selected edge attains this maximum, and $\xi_k\leq\theta G_{\mathcal A_k}(X^k)$ for a fixed $\theta<1$.
Then refinement terminates after at most
$$N_{\mathrm{ref}}\leq\left\lceil\frac{\Phi(X^0)-\Phi_{\inf}}{(1-\theta)\varepsilon_G}\right\rceil$$
rounds with $G_{\mathcal A}(X)\leq\varepsilon_G$.
\end{corollary}

\begin{proof}
Each nonterminal round decreases the objective by at least $(1-\theta)\varepsilon_G$.
The result follows by telescoping and the lower bound on $\Phi$.
\end{proof}

This alternative removes the uniform curvature assumption but terminates with a curvature-scaled certificate rather than the standard maximum dual violation.
If $h_e(X)\leq H$ on all inactive edges, then $G_{\mathcal A}(X)\leq\varepsilon_G$ implies $V_{\mathcal A}(X)\leq\sqrt{2H\varepsilon_G}$.

\section{Can PDUOT realize the inexact oracle?}

Fix one node-covering active support and consider the exact fixed-metric PDUOT core.
Its reflected Halpern iterates converge to a restricted KKT point, and the restricted optimal transported marginals are positive.
Consequently, the primal objective and the restricted projected-gradient residual converge to their optimal values.

\begin{proposition}[Oracle realizability]\label{prop:oracle-realizability}
Let $Y$ be any feasible coupling on a fixed node-covering active support.
For every $\varepsilon_r>0$ and $\xi>0$, an exact fixed-metric PDUOT run eventually produces a primal iterate $X$ satisfying
$$R_{\mathcal A}^{\gamma}(X)\leq\varepsilon_r,\qquad \Phi(X)\leq\Phi(Y)+\xi.$$
\end{proposition}

\begin{proof}
Let $X_{\mathcal A}^*$ be a restricted optimum.
Convergence of the PDUOT iterates gives $X^k\to X_{\mathcal A}^*$.
Positivity of the optimal marginals makes both $\Phi$ and $R_{\mathcal A}^{\gamma}$ continuous in a neighborhood of the limit, hence $R_{\mathcal A}^{\gamma}(X^k)\to0$ and $\Phi(X^k)\to\Phi(X_{\mathcal A}^*)$.
Because $Y$ is feasible, $\Phi(X_{\mathcal A}^*)\leq\Phi(Y)$.
Both required inequalities therefore hold after finitely many inner iterations.
\end{proof}

The proposition does not give an explicit inner-iteration bound because converting the fixed-point residual to primal objective error globally requires additional conditioning.
It nevertheless shows that the objective safeguard in Theorem~\ref{thm:inexact-refinement} is compatible with the convergent PDUOT core.
The inner tolerance can be tightened until both the restricted KKT condition and the round-dependent objective safeguard pass.

\section{Gap to the production refinement loop}

Theorem~\ref{thm:inexact-refinement} is not yet a theorem for the current HierUOT implementation unchanged.
The main difference is that PDUOT termination uses KKT measures and a split objective gap but does not enforce the round-dependent primal objective safeguard displayed in the refinement scheme.
The production loop therefore accepts the next restricted solution without checking that it retains a fixed fraction of the certified one-edge decrease.

Three smaller differences remain.
First, the production full-space scan uses PDUOT potentials, whereas the clean theorem uses primal-consistent potentials or the two-sided corrections in Corollary~\ref{cor:dual-correction}.
Second, production pruning protects flows only above a numerical threshold, whereas the proof preserves the exact current positive support and the selected descent edge.
Third, production stopping uses a relative $\ell_2$ residual by default and a relative $\ell_\infty$ residual as an option, whereas the theorem is stated with an absolute $\ell_\infty$ residual.
The norm and normalization difference changes only the constants on a finite-dimensional problem, but a floating-point theorem would also need an explicit full-scan error allowance.

Zero initialization of newly inserted edges is not a theoretical difference.
The analytic trial $Y^k$ only needs to be feasible on the expanded support, and the restricted solver need not use it as its warm start.
Adding multiple candidates is also compatible with the proof as long as the expanded support contains at least one certified descent edge.
The permanently protected structural edges already make the production supports node covering.

A certified implementation variant would compute primal-consistent node potentials after every restricted solve, evaluate the active projected-gradient residual, and use these potentials for the full-space scan.
When refinement is required, it would compute the analytic one-edge trial value without changing the warm start.
It would preserve the trial support through pruning and continue the restricted solve until both displayed oracle conditions are satisfied.

These changes are not needed to obtain the practical production result.
They define a controlled variant whose finite-tolerance guarantee is substantially closer to floating-point computation than an exact-restricted-oracle theorem.

\section{Assessment}

The one-edge descent lemma and Theorem~\ref{thm:inexact-refinement} provide a viable UOT-specific route to approximate-oracle convergence.
The key advantage over balanced OT is structural rather than technical: the soft KL marginals make every nonnegative coordinate direction feasible.
The price is a curvature constant governed by small transported marginals and an objective safeguard absent from the current code.

The remaining mathematical task is to replace the conservative global constant by a local or path-dependent complexity estimate that reflects the observed marginals.
The next implementation task is a diagnostic that measures whether production restricted solves already satisfy the round-dependent objective safeguard and how much positive flow is removed by threshold pruning.

\begin{thebibliography}{9}

\bibitem{liu2025etm}
W.~Liu, X.~Liao, J.~Dan, F.~Wang, H.~Yu, J.~Dong, S.~Dong, L.~Qi, and Y.-S.~Ong.
\newblock Solving Discrete (Semi) Unbalanced Optimal Transport with Equivalent Transformation Mechanism and KKT-Multiplier Regularization.
\newblock In \emph{Advances in Neural Information Processing Systems}, 2025.

\end{thebibliography}

\end{document}
